% Vietnamese translation for OI Public Library
% Source-PDF: IOI中国国家候选队论文/国家集训队2016论文集.pdf
\documentclass[11pt]{article}
\usepackage{oipl-vn}

\title{Tập luận văn đội tuyển quốc gia Trung Quốc năm 2016}
\author{Huấn luyện viên: Yu Linyun, Chen Xumin}
\date{China Computer Federation, 2016}

\begin{document}
\maketitle

\origtitle{Tập luận văn ứng viên đội tuyển quốc gia Trung Quốc Olympic Tin học năm 2016}
\sourcepath{IOI China candidate team papers / National training team 2016 collection.pdf}

\renewcommand{\abstractname}{Tóm tắt}
\begin{abstract}
PDF nguồn là tập hợp luận văn đội tuyển quốc gia Trung Quốc năm 2016. Phần
bìa và mục lục có lớp văn bản đọc được, nhưng ngay từ bài đầu tiên phần thân
đã bị lỗi ánh xạ phông chữ và trích xuất thành mojibake. Bản dịch này chuyển
ngữ các thông tin đầu sách, mục lục và bài đầu tiên của Ren Zhizhou về một số
phương pháp tính tổng hàm nhân tính; phần thân bài được dựng từ OCR và đối
chiếu công thức với lớp văn bản PDF.
\end{abstract}

\section{Thông tin đầu sách}

Tập sách có tiêu đề \emph{Tập luận văn ứng viên đội tuyển quốc gia Trung Quốc
Olympic Tin học năm 2016}. Huấn luyện viên ghi trên bìa là Yu Linyun và Chen
Xumin. Đơn vị xuất bản là China Computer Federation.

\section{Mục lục đã dịch}

\begin{center}
\small
\begin{tabular}{r p{0.47\linewidth} p{0.28\linewidth}}
\textbf{Trang} & \textbf{Bài viết} & \textbf{Tác giả / trường} \\
\hline
1 & Một số phương pháp tính tổng hàm nhân tính & Ren Zhizhou, Trường Trung học số 1 Thiệu Hưng \\
17 & Một số phương pháp mô hình hóa bằng luồng mạng & Jiang Zhihao, Trường Trung học số 1 Shengli, Đông Doanh \\
33 & Bàn sơ lược về quy hoạch tuyến tính và bài toán đối ngẫu & Dong Kefan, Trường Trung học số 1 Phúc Châu, Phúc Kiến \\
63 & Bàn sơ lược về một số thuật toán và ứng dụng cho bài toán cắt nhỏ nhất trên đồ thị vô hướng & Wang Wentao, Trường Trung học số 1 Thiệu Hưng \\
81 & Bàn sơ lược về ứng dụng của quy hoạch tuyến tính trong thi tin học & Zou Xiaoyao, Trường Trung học Trấn Hải, Ninh Ba \\
99 & Phép toán cực trị trên đoạn và bài toán cực trị lịch sử & Ji Ruyi, Trường Học Quân Hàng Châu \\
123 & Xét lại biến đổi Fourier nhanh & Mao Xiao, Trường Yali \\
141 & Từ \emph{Unknown} bàn về một lớp phương pháp duy trì thông tin đoạn có hỗ trợ chèn và xóa ở cuối & Luo Zhezheng, Trường THPT trực thuộc Đại học Sư phạm An Huy \\
163 & Báo cáo ra đề mảng hậu tố của Tiểu C & Hong Huadun, Trường Trung học số 1 Thiệu Hưng \\
171 & Báo cáo ra đề \emph{Xiaoxiaokan} & Zhang Haowei, Trường Trung học Dư Diêu, Chiết Giang \\
179 & Báo cáo ra đề \emph{strakf} & Li Zihao, Trường Trung học Shimen, Phật Sơn \\
189 & Báo cáo ra đề \emph{Tập hợp của quá khứ} & Wang Wenxiao, Trường Trung học Song Thập, Hạ Môn \\
195 & Báo cáo ra đề \emph{Lái tàu rời Qin Chuan} & Wu Zuofan, Trường THPT trực thuộc Đại học Sư phạm An Huy \\
209 & Bài luyện tập thuật toán sắp xếp cơ bản & Jin Ce, Trường Học Quân Hàng Châu, Chiết Giang \\
221 & Báo cáo ra đề \emph{move} & Yuan Yutao, Trường Yali, Trường Sa
\end{tabular}
\end{center}

\section*{Bài 1. Một số phương pháp tính tổng hàm nhân tính}
\addcontentsline{toc}{section}{Bài 1. Một số phương pháp tính tổng hàm nhân tính}

\begin{center}
Ren Zhizhou, Trường Trung học số 1 Thiệu Hưng
\end{center}

\subsection*{Tóm tắt}

Hàm nhân tính là một lớp hàm số học đặc biệt. Bài viết này tổng hợp một số kỹ
thuật tính toán thường gặp trong các bài toán liên quan tới hàm nhân tính, và
từ cảm hứng của các thuật toán đó đưa ra một phương pháp tính tổng hàm nhân
tính có phạm vi áp dụng tương đối rộng.

\subsection*{1. Dẫn nhập}

Trong thi tin học, đôi khi ta gặp các hàm được định nghĩa trên miền số nguyên
dương, tức các hàm số học. Hàm nhân tính, với tư cách một lớp hàm số học đặc
biệt, thường là đối tượng được thảo luận. Trong loại bài toán này, thí sinh
thường phải trước hết thực hiện một số suy luận, rồi chọn phương pháp thích
hợp để tính biểu thức sau khi đã biến đổi. Khi đó ta cần những thuật toán hiệu
quả để hoàn tất bước tính toán cuối cùng.

Trong bài viết này, tác giả trước hết chỉnh lý một số phương pháp tính toán
thường dùng trong thi tin học.

Mục 2 giới thiệu định nghĩa hàm nhân tính và liệt kê một số ví dụ thường gặp.

Mục 3 giới thiệu cách dùng sàng tuyến tính để nhanh chóng tính các giá trị
đầu của một hàm nhân tính.

Mục 4 giới thiệu định nghĩa tích chập Dirichlet và đảo Mobius, làm nền cho
mục sau.

Mục 5 kết hợp nội dung hai mục trước để đưa ra một phương pháp tính tổng hàm
số học hiệu quả. Trong thi tin học Trung Quốc, thuật toán này thường được gọi
là sàng Du Jiao; nó xây dựng công thức bằng tích chập Dirichlet. Tuy không có
liên hệ quá lớn với tính nhân tính, và đòi hỏi chính hàm cần tính có tính chất
đặc biệt nên điều kiện sử dụng hơi nghiêm ngặt, nó vẫn đem lại nhiều gợi ý.

Trong Mục 6, từ cảm hứng của các thuật toán trước đó, tác giả đưa ra một
phương pháp tính tổng hàm nhân tính mới. Hiệu suất của phương pháp này tuy hơi
kém thuật toán ở Mục 5, nhưng phạm vi áp dụng rộng hơn. Ở mục này, tác giả bắt
đầu từ động cơ thiết kế thuật toán, từng bước xây dựng toàn bộ quy trình, rồi
giới thiệu ứng dụng của thuật toán này ngoài bài toán tính tổng hàm nhân tính.

\subsection*{2. Định nghĩa}

\paragraph{\term{Định nghĩa 2.1}.}
Hàm có miền xác định là tập số nguyên dương và miền giá trị là tập số phức
được gọi là \term{hàm số học}.

\paragraph{\term{Định nghĩa 2.2}.}
Giả sử \(f(x)\) là một hàm số học. Nếu với mọi cặp số nguyên dương nguyên tố
cùng nhau \(a,b\) đều có
\[
  f(ab)=f(a)f(b),
\]
thì \(f(x)\) là một \term{hàm nhân tính}.

\paragraph{\term{Định nghĩa 2.3}.}
Giả sử \(f(x)\) là một hàm số học. Nếu với mọi cặp số nguyên dương \(a,b\) đều
có
\[
  f(ab)=f(a)f(b),
\]
thì \(f(x)\) là một \term{hàm hoàn toàn nhân tính}.

Nếu không nói riêng, mọi hàm xuất hiện trong bài viết này đều là hàm số học.

\subsubsection*{2.1. Một số hàm nhân tính thường gặp}

Sau đây là một số hàm nhân tính thường gặp:
\begin{itemize}
  \item Hàm ước \(\sigma_k(n)\), biểu thị tổng lũy thừa bậc \(k\) của mọi ước
  của \(n\).
  \item Hàm Euler \(\varphi(n)\), biểu thị số lượng số nguyên dương không vượt
  quá \(n\) và nguyên tố cùng nhau với \(n\).
  \item Hàm Mobius
  \[
    \mu(n)=
    \begin{cases}
      1, & n=1,\\
      (-1)^k, & n\text{ là tích của }k\text{ số nguyên tố khác nhau},\\
      0, & \text{các trường hợp khác}.
    \end{cases}
  \]
\end{itemize}

\subsubsection*{2.2. Một số hàm hoàn toàn nhân tính thường gặp}

Sau đây là một số hàm hoàn toàn nhân tính thường gặp:
\begin{itemize}
  \item Hàm lũy thừa \(Id_k(n)=n^k\).
  \item Hàm đơn vị
  \[
    \epsilon(n)=
    \begin{cases}
      1, & n=1,\\
      0, & n>1.
    \end{cases}
  \]
\end{itemize}

\subsection*{3. Sàng tuyến tính}

Sàng tuyến tính là thuật toán rất thường dùng khi tính hàm nhân tính. Nó có
thể tính trong \(O(n)\) thừa số nguyên tố nhỏ nhất của mọi số từ \(2\) tới
\(n\); dùng các thông tin này, ta có thể truy hồi xong các giá trị hàm của
\(n\) hạng đầu.

\subsubsection*{3.1. Tính thừa số nguyên tố nhỏ nhất}

Xét phân tích thừa số nguyên tố của một số nguyên dương \(n\), và sắp các thừa
số nguyên tố theo thứ tự giảm dần:
\[
  n=\prod_{i=1}^{k} p_i,
\]
trong đó \(p_i\) là số nguyên tố và với \(1\le i<k\) có \(p_i\ge p_{i+1}\).
Cách phân tích này là duy nhất; có thể dùng tính duy nhất đó để thu được thuật
toán sau:
\begin{itemize}
  \item Gọi \(pr_i\) là thừa số nguyên tố nhỏ nhất của \(i\), và liệt kê \(i\)
  từ \(2\) để tính.
  \item Nếu \(pr_i\) của số hiện tại chưa được tính, thì \(i\) là số nguyên
  tố và \(pr_i=i\).
  \item Liệt kê mọi số nguyên tố \(p\) không vượt quá \(pr_i\), rồi gán
  \(pr_{ip}=p\).
\end{itemize}

Mỗi số chỉ bị sàng bởi thừa số nguyên tố nhỏ nhất của nó, nên độ phức tạp thời
gian của thuật toán này là \(O(n)\).

\subsubsection*{3.2. Tính \(n\) hạng đầu của một hàm nhân tính}

Giả sử \(f(x)\) là một hàm nhân tính, nay cần tính các giá trị \(f(x)\) của
\(n\) hạng đầu.

Nếu có thể phân tích \(x\) thành một cặp \(a,b\) nguyên tố cùng nhau với
\(a,b>1\), thì giá trị \(f(x)\) có thể truy hồi theo Định nghĩa 2.2. Vì vậy
chỉ còn cần tính trường hợp \(x\) chỉ có một loại thừa số nguyên tố.

Ta có thể dùng sàng tuyến tính \(O(n)\) để tính thừa số nguyên tố nhỏ nhất của
mỗi số, rồi truy hồi số mũ của thừa số nguyên tố nhỏ nhất đó trong mỗi số.
Giả sử thừa số nguyên tố nhỏ nhất của \(x\) là \(p\), số mũ là \(c\), và
\(p^c\mid x\). Khi đó
\[
  f(x)=f(p^c)f\left(\frac{x}{p^c}\right).
\]

\subsubsection*{3.3. Ví dụ 1}

\paragraph{\term{Ví dụ 1}.}
Tính \(n\) giá trị đầu của hàm Euler \(\varphi(n)\).

Với \(n>1\), giả sử
\[
  n=\prod_{i=1}^{k} p_i^{c_i},
\]
trong đó các \(p_i\) đều là số nguyên tố và \(p_i\ne p_j\) khi \(i\ne j\).
Hàm Euler có công thức kinh điển
\[
  \varphi(n)
  = n\prod_{i=1}^{k}\frac{p_i-1}{p_i}
  = \prod_{i=1}^{k} (p_i-1)p_i^{c_i-1}.
\]

Với trường hợp \(n\) chỉ có một loại thừa số nguyên tố, ta có thể tính trực
tiếp; phần còn lại dùng cách ở Mục 3.2 để truy hồi bằng sàng tuyến tính. Độ
phức tạp thời gian là \(O(n)\).

\subsection*{4. Đảo Mobius}

\subsubsection*{4.1. Tích chập Dirichlet}

\paragraph{\term{Định nghĩa 4.1}.}
Định nghĩa tích chập Dirichlet của hai hàm số học \(f(x),g(x)\), ký hiệu
\((f*g)(x)\), là
\[
  (f*g)(n)=\sum_{d\mid n} f(d)g\left(\frac{n}{d}\right).
\]

Tích chập Dirichlet có các tính chất phép toán sau:
\begin{itemize}
  \item Giao hoán: \(f*g=g*f\).
  \item Kết hợp: \((f*g)*h=f*(g*h)\).
  \item Phân phối: \(f*(g+h)=f*g+f*h\).
  \item Đơn vị: \(f*\epsilon=f\).
  \item Nếu \(f,g\) đều là hàm nhân tính thì \(f*g\) cũng là hàm nhân tính.
\end{itemize}

Các tính chất trên chứng minh khá đơn giản, xin để độc giả tự suy nghĩ.

\subsubsection*{4.2. Đảo Mobius}

\paragraph{\term{Bổ đề 4.1}.}
\(\mu*1=\epsilon\), tức
\[
  \sum_{d\mid n}\mu(d)=\epsilon(n),
\]
trong đó hàm \(1\) là hàm hằng luôn trả về \(1\).

\paragraph{\term{Chứng minh}.}
Giả sử \(n\) có \(k\) loại thừa số nguyên tố khác nhau. Khi đó
\[
  \sum_{d\mid n}\mu(d)
  = \sum_{i=0}^{k}(-1)^i\binom{k}{i}
  = (1-1)^k
  = [k=0]
  = \epsilon(n).
\]
\(\epsilon\) là hàm đơn vị đã nhắc tới ở Mục 2.2.

\paragraph{\term{Định lý 4.1} (đảo Mobius).}
Nếu có hai hàm số học \(f,g\) thỏa
\[
  f(n)=\sum_{d\mid n}g(d),
\]
thì chúng cũng thỏa
\[
  g(n)=\sum_{d\mid n}\mu(d)f\left(\frac{n}{d}\right).
\]
Điều ngược lại cũng đúng, tức
\[
  f=g*1 \Longleftrightarrow g=\mu*f.
\]

\paragraph{\term{Chứng minh}.}
Giả sử
\[
  f=g*1.
\]
Chập hai vế với hàm \(\mu\), ta được
\[
  \mu*f=\mu*g*1.
\]
Theo Bổ đề 4.1, suy ra
\[
  \mu*f=\epsilon*g=g.
\]
Ngược lại, chập hai vế với hàm \(1\), ta được
\[
  g*1=\mu*f*1=\epsilon*f=f.
\]

\subsubsection*{4.3. Mở rộng của đảo Mobius}

Đảo Mobius còn có nhiều dạng mở rộng. Tiếp theo ta giới thiệu một mở rộng về
hàm số học.

\paragraph{\term{Bổ đề 4.2}.}
Giả sử \(x,y,a,b\) là các số nguyên dương, trong đó \(a,b\le x\), và
\[
  y=\left\lfloor\frac{x}{a}\right\rfloor.
\]
Khi đó
\[
  \left\lfloor\frac{y}{b}\right\rfloor
  = \left\lfloor\frac{x}{ab}\right\rfloor.
\]

\paragraph{\term{Chứng minh}.}
Đặt \(x=kab+c\), trong đó \(k,c\) là các số nguyên không âm và \(c<ab\). Khi
đó \(k=\left\lfloor x/(ab)\right\rfloor\). Ta có
\[
  y=\left\lfloor\frac{x}{a}\right\rfloor
  = kb+\left\lfloor\frac{c}{a}\right\rfloor,
\]
và \(\left\lfloor c/a\right\rfloor<b\), nên
\(\left\lfloor y/b\right\rfloor=k\).

\paragraph{\term{Định lý 4.2}.}
Giả sử \(f,g\) là hai hàm số học, \(t\) là một hàm hoàn toàn nhân tính và
\(t(1)=1\). Khi đó
\[
  f(n)=\sum_{k=1}^{n} t(k)g\left(\left\lfloor\frac{n}{k}\right\rfloor\right)
\]
khi và chỉ khi
\[
  g(n)=\sum_{k=1}^{n}\mu(k)t(k)
  f\left(\left\lfloor\frac{n}{k}\right\rfloor\right).
\]

\paragraph{\term{Chứng minh}.}
Xét thay công thức đầu vào công thức sau. Theo Bổ đề 4.2, ta có
\[
\begin{aligned}
  \sum_{k=1}^{n}\mu(k)t(k)
  f\left(\left\lfloor\frac{n}{k}\right\rfloor\right)
  &=
  \sum_{k=1}^{n}\mu(k)t(k)
  \sum_{i=1}^{\lfloor n/k\rfloor}
  t(i)g\left(\left\lfloor\frac{n}{ki}\right\rfloor\right)\\
  &=
  \sum_{j=1}^{n}
  g\left(\left\lfloor\frac{n}{j}\right\rfloor\right)
  \sum_{i\mid j}\mu(i)t(i)t\left(\frac{j}{i}\right)\\
  &=
  \sum_{j=1}^{n}
  g\left(\left\lfloor\frac{n}{j}\right\rfloor\right)
  t(j)\sum_{i\mid j}\mu(i).
\end{aligned}
\]
Theo Bổ đề 4.1,
\[
  \sum_{k=1}^{n}\mu(k)t(k)
  f\left(\left\lfloor\frac{n}{k}\right\rfloor\right)
  =
  \sum_{j=1}^{n}
  g\left(\left\lfloor\frac{n}{j}\right\rfloor\right)t(j)\epsilon(j)
  =g(n)t(1)=g(n).
\]
Chiều ngược lại cũng tương tự.

\subsection*{5. Xây dựng bằng tích chập Dirichlet}

Trong thi tin học Trung Quốc, thuật toán này thường được gọi là sàng Du Jiao.
Nó có thể hoàn thành hiệu quả phép tính cho một số hàm số học thường gặp.

\subsubsection*{5.1. Dạng chính}

Giả sử \(f(n)\) là một hàm số học, cần tính
\[
  S(n)=\sum_{i=1}^{n} f(i).
\]
Theo tính chất của hàm \(f(n)\), ta xây dựng một công thức truy hồi của
\(S(n)\) theo các giá trị
\[
  S\left(\left\lfloor\frac{n}{i}\right\rfloor\right),
\]
như dưới đây.

Tìm một hàm số học thích hợp \(g(n)\):
\[
  \sum_{i=1}^{n}\sum_{d\mid i} f(d)g\left(\frac{i}{d}\right)
  =
  \sum_{i=1}^{n}g(i)S\left(\left\lfloor\frac{n}{i}\right\rfloor\right).
\]
Khi đó thu được công thức truy hồi
\[
  g(1)S(n)
  =
  \sum_{i=1}^{n}(f*g)(i)
  -
  \sum_{i=2}^{n}g(i)S\left(\left\lfloor\frac{n}{i}\right\rfloor\right).
\]

\paragraph{\term{Bổ đề 5.1}.}
Với một số nguyên dương \(n\), xét mọi số nguyên dương \(i\in[1,n]\). Các giá
trị \(\left\lfloor n/i\right\rfloor\) chỉ có \(O(\sqrt n)\) loại.

\paragraph{\term{Chứng minh}.}
Với mọi \(i\le\sqrt n\), ta có \(\left\lfloor n/i\right\rfloor\ge\sqrt n\),
nên phần này có \(O(\sqrt n)\) loại giá trị. Với mọi \(i>\sqrt n\), ta có
\(\left\lfloor n/i\right\rfloor<\sqrt n\), nên phần này cũng có
\(O(\sqrt n)\) loại giá trị.

Theo Bổ đề 4.2 và Bổ đề 5.1, trong toàn bộ quá trình truy hồi tính \(S(n)\),
chỉ có \(O(\sqrt n)\) giá trị
\[
  S\left(\left\lfloor\frac{n}{i}\right\rfloor\right)
\]
cần được tính.

Nếu có thể nhanh chóng lấy tổng tiền tố của \((f*g)(i)\) và \(g(i)\), ta có
thể chia đoạn theo giá trị của \(\left\lfloor n/i\right\rfloor\). Khi đó tính
một \(S(n)\) có độ phức tạp \(O(\sqrt n)\), còn toàn bộ quá trình có độ phức
tạp
\[
  \sum_{i=1}^{\lfloor\sqrt n\rfloor} O(\sqrt i)
  + \sum_{i=1}^{\lfloor\sqrt n\rfloor}
  O\left(\sqrt{\frac{n}{i}}\right).
\]
Phần sau chắc chắn lớn hơn phần trước, nên chỉ cần xét
\[
  \sum_{i=1}^{\lfloor\sqrt n\rfloor}
  O\left(\sqrt{\frac{n}{i}}\right)
  \approx
  O\left(\int_0^{\sqrt n}\sqrt{\frac{n}{x}}\,dx\right)
  = O(n^{3/4}).
\]

Cách xây dựng công thức truy hồi cần phân tích theo hàm cụ thể. Có thể tham
khảo các ví dụ sau.

\subsubsection*{5.2. Ví dụ 2}

\paragraph{\term{Ví dụ 2}.}
Tính
\[
  S(n)=\sum_{i=1}^{n}\mu(i).
\]

Theo Bổ đề 4.1, dùng phương pháp ở Mục 5.1 ta thu được công thức truy hồi
\[
  S(n)=\sum_{i=1}^{n}\epsilon(i)
  -\sum_{i=2}^{n}S\left(\left\lfloor\frac{n}{i}\right\rfloor\right)
  =
  1-\sum_{i=2}^{n}S\left(\left\lfloor\frac{n}{i}\right\rfloor\right).
\]
Độ phức tạp tính trực tiếp là \(O(n^{3/4})\). Xét việc dùng sàng tuyến tính xử
lý trước \(n^{1/3}\) hạng đầu, khi đó độ phức tạp của phần truy hồi là
\[
  O\left(\int_0^{n^{1/3}}\sqrt{\frac{n}{x}}\,dx\right)=O(n^{2/3}).
\]
Kết hợp với phần sàng tuyến tính, tổng độ phức tạp là \(O(n^{2/3})\).

\subsubsection*{5.3. Ví dụ 3}

\paragraph{\term{Ví dụ 3} (SPOJ DIVCNT2).}
Tính
\[
  S(n)=\sum_{i=1}^{n}\sigma_0(i^2).
\]

Tương tự, xét xây dựng quan hệ tích chập Dirichlet. Đặt
\[
  f(n)=\sigma_0(n^2),\qquad g(n)=2^{\omega(n)},
\]
trong đó \(\omega(n)\) biểu thị số lượng thừa số nguyên tố khác nhau của
\(n\). Ta có \(f=g*1\), tức
\[
  \sigma_0(n^2)=\sum_{d\mid n}2^{\omega(d)}.
\]

\paragraph{\term{Chứng minh}.}
Xét từng ước \(d\) của \(n\). Trong \(d^2\), ta có thể bỏ đi một tập thừa số
nguyên tố của \(d\), tức có \(2^{\omega(d)}\) cách bỏ; bằng cách này có thể
liệt kê mọi ước của \(n^2\).

Dễ thấy \(g=\mu^2*1\), do đó ta có các quan hệ truy hồi sau:
\[
\begin{aligned}
  \sum_{i=1}^{n}\sigma_0(i^2)
  &= \sum_{i=1}^{n}2^{\omega(i)}
     \left\lfloor\frac{n}{i}\right\rfloor,\\
  \sum_{i=1}^{n}2^{\omega(i)}
  &= \sum_{i=1}^{n}\mu^2(i)
     \left\lfloor\frac{n}{i}\right\rfloor,\\
  \sum_{i=1}^{n}\mu^2(i)
  &= \sum_{i=1}^{\lfloor\sqrt n\rfloor}\mu(i)
     \left\lfloor\frac{n}{i^2}\right\rfloor.
\end{aligned}
\]
Tương tự Ví dụ 2, có thể dùng sàng tuyến tính để hoàn thành tính toán trong
\(O(n^{2/3})\).\footnote{Công thức tính tổng của \(\mu^2\) có thể nhận được
từ bao hàm--loại trừ.}

\subsubsection*{5.4. Tiểu kết}

Theo Bổ đề 4.2 và Bổ đề 5.1, khi truy hồi tính hạng thứ \(n\), loại công thức
này chỉ liên quan tới \(O(\sqrt n)\) trạng thái. Bằng tích phân, có thể tính
độ phức tạp của chuyển trạng thái truy hồi là \(O(n^{3/4})\); dùng sàng tuyến
tính xử lý trước một đoạn đầu có thể tiếp tục hạ độ phức tạp.

Lấy tư tưởng này làm nền tảng, ta có thể xây dựng thuật toán trong mục sau.
Về nội dung của mục này, có thể xem tài liệu tham khảo [2] để biết thêm chi
tiết; bài viết này không triển khai sâu hơn.

\subsection*{6. Một phương pháp tính tổng dựa trên phân tích thừa số nguyên tố}

Trong nhiều trường hợp, ta không thể dễ dàng tìm được tính chất của hàm cần
tính tổng. Khi đó ta hy vọng có một phương pháp tính tổng hiệu quả với phạm vi
áp dụng rộng hơn. Tiếp theo bài viết đưa ra một thuật toán tính tổng cho các
hàm nhân tính có dạng tương đối tổng quát.

\subsubsection*{6.1. Nhắc lại}

Giả sử \(F(n)\) là một hàm nhân tính, và phân tích thừa số nguyên tố của \(n\)
là
\[
  n=\prod_{i=1}^{k} p_i^{c_i}.
\]
Xét cách mô tả một hàm nhân tính như sau:
\begin{itemize}
  \item Khi \(n\) là số nguyên tố, tức \(n=p\), \(F(p)=G(p)\).
  \item Khi \(n\) là lũy thừa của một số nguyên tố, tức \(n=p^c\) với
  \(c>1\), \(F(p^c)=T(p^c)\).
  \item Các trường hợp còn lại tính theo tính nhân tính:
  \[
    F(n)=\prod_{i=1}^{k}F(p_i^{c_i}).
  \]
\end{itemize}

Ví dụ, với hai hàm nhân tính thường gặp sau:
\begin{itemize}
  \item Với hàm Euler \(\varphi(n)\), \(G(p)=p-1\) và
  \(T(p^c)=(p-1)p^{c-1}\).
  \item Với hàm Mobius \(\mu(n)\), \(G(p)=-1\) và \(T(p^c)=0\).
\end{itemize}

Tổng quát hơn, \(G(p)\) có thể là một đa thức theo \(p\), còn \(T(p^c)\) có
thể là một đa thức theo \(p\) và \(c\). Với các trường hợp này, hoặc với những
biểu thức phức tạp hơn, thuật toán ở mục trước có thể không còn phù hợp. Ta
xét không dựa vào tính chất riêng của hàm nữa, mà trực tiếp hoàn thành tính
toán bằng tính nhân tính.

\subsubsection*{6.2. Đưa vào bài toán}

Lấy một dạng đơn giản của loại bài toán này làm ví dụ.

\paragraph{\term{Ví dụ 4}.}
Định nghĩa một biến thể \(\phi(n,d)\) của hàm Euler. Phân tích thừa số nguyên
tố của \(n\) là
\[
  n=\prod_{i=1}^{k}p_i^{c_i},
\]
trong đó \(p_i\) là các thừa số nguyên tố đôi một khác nhau và \(c_i>0\). Với
trường hợp \(n=1\), ta hiểu tích sau là tích rỗng. Đặt
\[
  \phi(n,d)=\prod_{i=1}^{k}(p_i^{c_i}+d).
\]
Đặc biệt, định nghĩa \(\phi(1,d)=1\). Với \(n,d\) cho trước, hãy tính
\[
  \left(\sum_{i=1}^{n}\phi(i,d)\right)\bmod (10^9+7).
\]

Đặt \(F(n)=\phi(n,d)\). Theo cách biểu diễn ở Mục 6.1, ta có
\[
  G(p)=p+d,\qquad T(p^c)=p^c+d.
\]

\subsubsection*{6.3. Biến đổi bước đầu}

\paragraph{\term{Bổ đề 6.1}.}
Với một số nguyên dương \(x\le n\), \(x\) có nhiều nhất một thừa số nguyên tố
lớn hơn \(\sqrt n\).

\paragraph{\term{Chứng minh}.}
Giả sử \(x\) có hai thừa số nguyên tố \(p_1,p_2\) lớn hơn \(\sqrt n\). Khi đó
\(p_1p_2>n\), mâu thuẫn với \(x\le n\).

Theo Bổ đề 6.1, có thể chia \(F(x)\) thành hai loại để xét:
\begin{itemize}
  \item \(x\) không có thừa số nguyên tố lớn hơn \(\sqrt n\);
  \item \(x\) có một thừa số nguyên tố lớn hơn \(\sqrt n\).
\end{itemize}

Theo tính nhân tính của hàm, ta có
\[
  \sum_{i=1}^{n}F(i)
  =
  \sum_{\substack{x\le n\\x\text{ không có thừa số nguyên tố }>\sqrt n}}
  F(x)
  \left(
    1+\sum_{\substack{\sqrt n<p\le \lfloor n/x\rfloor\\p\text{ nguyên tố}}}
    G(p)
  \right).
\]

Hệ số nhân với \(F(x)\) chỉ liên quan tới giá trị
\(\left\lfloor n/x\right\rfloor\). Vì vậy chỉ có \(O(\sqrt n)\) loại hệ số
khác nhau, và ta có thể phân \(F(x)\) thành \(O(\sqrt n)\) nhóm theo giá trị
\(\left\lfloor n/x\right\rfloor\), mỗi nhóm ứng với một loại hệ số.

Đặt \(y\) là một giá trị cố định của \(\left\lfloor n/x\right\rfloor\). Khi đó
cần tính hai phần
\[
  \sum_{\substack{\sqrt n<p\le y\\p\text{ nguyên tố}}}G(p),
  \qquad
  \sum_{\substack{\left\lfloor n/x\right\rfloor=y\\
  x\text{ không có thừa số nguyên tố }>\sqrt n}}F(x).
\]

\subsubsection*{6.4. Tính \(G(p)\)}

\paragraph{\term{6.4.1. Phân tích bước đầu}.}
Trong công thức của Ví dụ 4, \(G(p)=p+d\). Không mất tính tổng quát, giả sử
\(G(p)\) là một đa thức bậc thấp theo \(p\), và xét tính riêng từng hạng của
đa thức.

Giả sử các số nguyên tố không vượt quá \(\sqrt n\) có tổng cộng \(m\) số, sắp
theo thứ tự tăng dần là \(p_1,\ldots,p_m\). Đặt \(g_k[i][j]\) là tổng lũy thừa
bậc \(k\) của mọi số trong đoạn \([1,j]\) nguyên tố cùng nhau với \(i\) số
nguyên tố đầu. Phần \(i=0\) có thể tính bằng nội suy trong \(O(k)\).

Với \(i\ge1\), xét trong đoạn \([1,j]\) có bao nhiêu số chứa thừa số nguyên tố
\(p_i\) và nguyên tố cùng nhau với \(i-1\) số nguyên tố trước. Đặt
\[
  l=\left\lfloor\frac{j}{p_i}\right\rfloor.
\]
Ta có truy hồi
\[
  g_k[i][j]=g_k[i-1][j]-p_i^k g_k[i-1][l].
\]

Giá trị cuối cùng \(g_k[m][j]-1\) chính là tổng lũy thừa bậc \(k\) của các số
nguyên tố lớn hơn \(\sqrt n\) trong đoạn \([1,j]\). Theo Bổ đề 4.2 và Bổ đề
5.1, trong quá trình tính \(g_k[m][n]\), số trạng thái chiều thứ hai cần dùng
là \(O(\sqrt n)\), và các trạng thái này đúng là từng loại tổng \(G(p)\) cần
tính.

\paragraph{\term{6.4.2. Cài đặt thô}.}
Cài đặt thô công thức truy hồi này cần lần lượt liệt kê số nguyên tố \(p_i\)
và tính trên từng trạng thái.

Xét với mỗi giá trị \(\left\lfloor n/x\right\rfloor\) có bao nhiêu số nguyên
tố cần chuyển. Phần \(\left\lfloor n/x\right\rfloor>\sqrt n\) cần chuyển mọi
số nguyên tố trong phạm vi \(\sqrt n\), nên
\[
  \sum_{i=1}^{\lfloor\sqrt n\rfloor} O\left(\frac{i}{\log i}\right)
  +
  \sum_{i=1}^{\lfloor\sqrt n\rfloor}
  O\left(\frac{\sqrt n}{\log\sqrt n}\right)
  \approx
  O\left(\frac{n}{\log n}\right).
\]

\paragraph{\term{6.4.3. Tối ưu}.}
Khi \(p_{i+1}>j\), chắc chắn \(g_k[i][j]=1\). Vì vậy nếu
\(p_i^2>j\ge p_i\), tức
\[
  l=\left\lfloor\frac{j}{p_i}\right\rfloor<p_i,
\]
thì
\[
  g_k[i-1][l]=1,\qquad
  g_k[i][j]=g_k[i-1][j]-p_i^k.
\]

Do đó thực ra có thể bỏ qua tính toán khi \(p_i^2>j\). Với mỗi \(j\), ghi lại
giá trị \(i\) ở lần chuyển cuối cùng; khi gọi vị trí \(g_k[i][j]\), cộng bù
đóng góp của đoạn số nguyên tố chưa được tính.

Với mỗi giá trị \(\left\lfloor n/x\right\rfloor\), chỉ cần chuyển các số
nguyên tố không vượt quá \(\sqrt{\left\lfloor n/x\right\rfloor}\). Độ phức tạp
có thể ước lượng là
\[
  \sum_{i=1}^{\lfloor\sqrt n\rfloor} O\left(\frac{\sqrt i}{\log i}\right)
  +
  \sum_{i=1}^{\lfloor\sqrt n\rfloor}
  O\left(\frac{\sqrt{n/i}}{\log\sqrt{n/i}}\right).
\]
Đóng góp của phần sau chắc chắn lớn hơn phần trước, và \(\sqrt{n/i}\ge
n^{1/4}\), nên có thể ước lượng độ phức tạp bởi
\[
  \sum_{i=1}^{\lfloor\sqrt n\rfloor}
  O\left(\frac{\sqrt{n/i}}{\log\sqrt{n/i}}\right)
  \approx
  O\left(\frac{\int_0^{\sqrt n}\sqrt{n/x}\,dx}{\log n}\right)
  = O\left(\frac{n^{3/4}}{\log n}\right).
\]

\subsubsection*{6.5. Tính \(F(x)\)}

\paragraph{\term{6.5.1. Phân tích bước đầu}.}
Vì các \(F(x)\) có cùng giá trị \(\left\lfloor n/x\right\rfloor\) cần nhân
cùng một hệ số, có thể trực tiếp dùng \(\left\lfloor n/x\right\rfloor\) để
biểu diễn trạng thái.

Tương tự, giả sử các số nguyên tố không vượt quá \(\sqrt n\) có tổng cộng
\(m\) số, sắp theo thứ tự tăng dần là \(p_1,\ldots,p_m\). Đặt \(f[i][j]\) là
tổng các \(F(x)\) sao cho \(x\) chỉ chứa \(i\) loại thừa số nguyên tố đầu và
\(\left\lfloor n/x\right\rfloor=j\). Trong đó \(f[0][n]=1\).

Theo Bổ đề 4.2 và Bổ đề 5.1, số trạng thái chiều thứ hai vẫn có thể thu gọn
xuống \(O(\sqrt n)\); mỗi loại ứng với một nhóm hệ số.

\paragraph{\term{6.5.2. Cài đặt thô}.}
Xét chuyển từ \(f[i-1][j]\) bằng số nguyên tố \(p_i\). Liệt kê số mũ chuyển
\(c\), đặt
\[
  l=\left\lfloor\frac{j}{p_i^c}\right\rfloor.
\]
Khi đó cộng vào \(f[i][l]\) giá trị
\[
  T(p_i^c)f[i-1][j],
\]
hoặc \(G(p_i)f[i-1][j]\) trong trường hợp \(c=1\).

Tương tự như trước, tính xem mỗi giá trị \(\left\lfloor n/x\right\rfloor\)
có bao nhiêu số nguyên tố cần chuyển. Xét lượng tính toán khi liệt kê từng
loại \(c\), đặt
\[
  h(n)=\sum_{k=2}^{\lfloor\log_2 n\rfloor} O(n^{1/k})\approx O(\sqrt n).
\]
Ước lượng độ phức tạp tương tự Mục 6.4:
\[
  \sum_{i=1}^{\lfloor\sqrt n\rfloor}
  O\left(\frac{i+h(i)}{\log i}\right)
  +
  \sum_{i=1}^{\lfloor\sqrt n\rfloor}
  O\left(\frac{\sqrt n+h(n/i)}{\log n}\right)
  \approx
  O\left(\frac{n}{\log n}\right).
\]

\paragraph{\term{6.5.3. Tối ưu}.}
Ở tiểu mục trước, bằng cách bỏ qua tính toán khi \(p_i^2>j\), ta hạ độ phức
tạp xuống \(O(n^{3/4}/\log n)\). Xét áp dụng tối ưu tương tự.

Chú ý rằng khi \(\left\lfloor n/x\right\rfloor\) không vượt quá \(\sqrt n\),
\[
  1+\sum_{\substack{\sqrt n<p\le \lfloor n/x\rfloor\\p\text{ nguyên tố}}}G(p)
  =1.
\]
Mà công thức tính đáp án cuối cùng là
\[
  \sum_{i=1}^{n}F(i)
  =
  \sum_{\substack{x\le n\\x\text{ không có thừa số nguyên tố }>\sqrt n}}
  F(x)
  \left(
    1+\sum_{\substack{\sqrt n<p\le \lfloor n/x\rfloor\\p\text{ nguyên tố}}}
    G(p)
  \right).
\]
Hệ số đóng góp của phần \(F(x)\) này vào đáp án là giống nhau, nên có thể dùng
đó làm điểm đột phá để tối ưu.

Khi \(p_i^2>y=\left\lfloor n/x\right\rfloor\), ta có
\(\left\lfloor y/p_i\right\rfloor<p_i\le\sqrt n\), nên trạng thái \(y\) nhiều
nhất chỉ có thể chuyển bằng một số nguyên tố vượt quá \(p_i\), và giá trị sau
khi chuyển chắc chắn có hệ số đóng góp vào đáp án bằng \(1\).

Chiến lược chuyển sau khi tối ưu như sau:
\begin{itemize}
  \item Với một số nguyên tố \(p_i\), chỉ chuyển các trạng thái \(y\ge p_i^2\).
  \item Đặt \(l=\left\lfloor y/p_i^c\right\rfloor\). Nếu \(l<p_i^2\), thì
  trạng thái này sẽ bị bỏ qua trong các lần chuyển sau; cần tính ngay đóng góp
  của nó vào đáp án, tức thống kê tổng \(G(p)\) trong đoạn \([p_{i+1},l]\).
  \item Duy trì với mỗi trạng thái số nguyên tố \(p_i\) ở lần chuyển cuối
  cùng, rồi thống kê tổng \(G(p)\) của đoạn bị bỏ qua.
\end{itemize}

Độ phức tạp được ước lượng tương tự Mục 6.4:
\[
  \sum_{i=1}^{\lfloor\sqrt n\rfloor}O\left(\frac{h(i)}{\log i}\right)
  +
  \sum_{i=1}^{\lfloor\sqrt n\rfloor}
  O\left(\frac{h(n/i)}{\log n}\right)
  \approx
  O\left(\frac{n^{3/4}}{\log n}\right).
\]

\paragraph{\term{6.5.4. Một cách cài đặt khác}.}
Chú ý rằng chỉ khi \(x<\sqrt n\), ta mới có thể tiếp tục nhân thêm một số
nguyên tố lớn hơn \(\sqrt n\); phần \(x\) này cũng không thể chứa bất kỳ thừa
số nguyên tố nào lớn hơn \(\sqrt n\).

Có thể trước hết dùng sàng tuyến tính để xử lý giá trị \(F(x)\) của mọi
\(x<\sqrt n\), nhân hệ số tương ứng rồi cộng vào đáp án. Phần còn lại là tính
\[
  \sum_{\substack{\sqrt n\le x\le n\\
  x\text{ không có thừa số nguyên tố }>\sqrt n}}F(x).
\]

Khác với mục trước, sắp các số nguyên tố không vượt quá \(\sqrt n\) theo thứ
tự giảm dần là \(p_1,\ldots,p_m\). Đặt \(f'[i][j]\) là tổng các \(F(x)\) sao
cho chỉ chứa \(i\) loại thừa số nguyên tố đầu và \(x\le j\). Dễ có
\(f'[0][j]=1\); giá trị \(f'[m][n]\) sau khi trừ tổng các \(F(x)\) với
\(x<\sqrt n\) chính là kết quả cần tìm.

Dễ thu được một chuyển truy hồi: liệt kê số mũ \(c\) của thừa số nguyên tố
\(p_i\), đặt \(l_c=\left\lfloor j/p_i^c\right\rfloor\), ta có
\[
  f'[i][j]
  =
  f'[i-1][j]
  +G(p_i)f'[i-1][l_1]
  +\sum_{c\ge2}T(p_i^c)f'[i-1][l_c].
\]

Tham khảo mục trước, chiều thứ hai của truy hồi này cũng chỉ có \(O(\sqrt n)\)
trạng thái. Cài đặt thô có độ phức tạp \(O(n/\log n)\); xét áp dụng tối ưu
tương tự.

Khi \(p_i^2>j\), vì \(p_i\) được sắp theo thứ tự giảm dần, chỉ cần thống kê
tổng giá trị hàm của các số nguyên tố trong đoạn \([p_i,j]\), nên vẫn có thể
bỏ qua tính toán của phần này và hạ độ phức tạp xuống
\[
  O\left(\frac{n^{3/4}}{\log n}\right).
\]

So sánh hai cách cài đặt, khác biệt chính nằm ở thứ tự liệt kê thừa số nguyên
tố. Có thể chọn theo tình huống cụ thể.\footnote{Tác giả cảm ơn Mao Xiao đã
cung cấp cách cài đặt này.}

\subsubsection*{6.6. Mở rộng}

Nhắc lại vài bước biến đổi chính trong thuật toán này:
\begin{itemize}
  \item Theo Bổ đề 6.1, dùng tính nhân tính để chia thừa số nguyên tố của mỗi
  số thành hai phần để xét.
  \item Theo Bổ đề 4.2 và Bổ đề 5.1, thu nhỏ chiều thứ hai của truy hồi xuống
  \(O(\sqrt n)\) trạng thái.
  \item Phân tích tính chất của truy hồi, rồi bỏ qua tính toán khi \(p_i^2>j\)
  để hạ độ phức tạp xuống \(O(n^{3/4}/\log n)\).
\end{itemize}

Thực chất, các bước trên đã tách thừa số nguyên tố cho mọi số, và khi truy hồi
thì các thừa số nguyên tố được đưa vào có thứ tự. Vì vậy thuật toán này có thể
hoàn thành một số chức năng của sàng.

\paragraph{\term{6.6.1. Ví dụ 5}.}
\term{Ví dụ 5 (UR13C Sanrd).}

Tính tổng thừa số nguyên tố lớn thứ hai của mọi hợp số không vượt quá \(n\).
Ở đây thừa số nguyên tố lớn thứ hai được định nghĩa như sau: giả sử
\[
  n=\prod_{i=1}^{m}p_i,
\]
trong đó các \(p_i\) đều là số nguyên tố và \(p_i\le p_{i+1}\). Thừa số
nguyên tố lớn thứ hai là \(p_{m-1}\); nếu \(m<2\) thì không tính.

Tương tự, chia thừa số nguyên tố thành hai phần theo ngưỡng \(\sqrt n\), và
theo Bổ đề 6.1, thừa số nguyên tố lớn thứ hai của \(x\le n\) chắc chắn không
vượt quá \(\sqrt n\).

Nhắc lại thuật toán đầu tiên trong Mục 6.4 và Mục 6.5. Quy trình tạo thành
mỗi số \(x\) là: trước hết thêm các thừa số nguyên tố không vượt quá
\(\sqrt n\) của \(x\) theo thứ tự tăng dần, cuối cùng mới tính thừa số nguyên
tố lớn hơn \(\sqrt n\) của \(x\). Tức trong toàn bộ quá trình, các thừa số
nguyên tố của \(x\) được liệt kê theo thứ tự tăng dần.

Khi tính các thừa số nguyên tố không vượt quá \(\sqrt n\), có thể xem mỗi
trạng thái là đang duy trì thông tin của một tập số. Một lần chuyển là nhân
mọi số trong tập đó với thừa số nguyên tố đang được liệt kê, thu được một tập
mới và nhập vào trạng thái đích. Vì thừa số nguyên tố được liệt kê có thứ tự,
mỗi trạng thái sẽ không có số bị tính lặp, và còn bảo đảm các tập số giữa các
trạng thái không giao nhau.

Dễ thấy khi chuyển, thừa số nguyên tố vừa thêm chắc chắn là thừa số nguyên tố
lớn nhất của số mới sinh; thừa số nguyên tố lớn nhất của trạng thái được
chuyển chính là thừa số nguyên tố lớn thứ hai. Vì vậy chỉ cần ghi tổng thừa
số nguyên tố lớn nhất của tập số mà mỗi trạng thái biểu diễn, ta có thể thu
được thừa số nguyên tố lớn thứ hai.

\subsubsection*{6.7. Tổng kết}

Thuật toán này là nội dung chính bài viết muốn trình bày. Nó giải quyết bài
toán tính tổng hàm nhân tính với dạng tương đối tổng quát. Các ràng buộc chính
đối với hàm ban đầu như sau:
\begin{itemize}
  \item Hàm cần tính tổng \(F(n)\) phải là hàm nhân tính.
  \item \(G(p)\) có thể phân tích thành tổng của một số hàm hoàn toàn nhân
  tính; do bước khởi tạo mảng truy hồi cần điều này, tổng tiền tố của các hàm
  hoàn toàn nhân tính ấy phải tính nhanh được. Một ví dụ đơn giản là \(G(p)\)
  là đa thức bậc thấp theo \(p\).
  \item Do chuyển ở Mục 6.5 yêu cầu, hệ số chuyển \(G(p)\) và \(T(p^c)\) phải
  tính nhanh được.
\end{itemize}

Các chuyển truy hồi trong Mục 6.4 và Mục 6.5 đều có thể dùng mảng cuốn; số
nguyên tố cũng chỉ cần xử lý trước trong phạm vi \(\sqrt n\), nên độ phức tạp
không gian là \(O(\sqrt n)\).

Thuật toán này có thể giúp ta tính tổng đa số hàm nhân tính trong
\[
  O\left(\frac{n^{3/4}}{\log n}\right),
\]
và có ưu thế lớn với những bài toán tính tổng hàm nhân tính không có tính chất
đặc biệt hoặc khó suy luận.

\subsection*{Lời cảm ơn}

Xin cảm ơn Hiệp hội Máy tính đã cung cấp nền tảng học tập và giao lưu.

Xin cảm ơn các thầy Chen Heli và Dong Yehua của Trường Trung học số 1 Thiệu
Hưng vì sự quan tâm và chỉ dẫn trong nhiều năm.

Xin cảm ơn các anh Yu Yili, Dong Honghua, Zhang Hengjie và Wang Jianhao của
Đại học Thanh Hoa vì đã giúp đỡ tôi.

Xin cảm ơn bạn Mao Xiao đã thảo luận thuật toán này với tôi và cung cấp một
cách cài đặt ngắn gọn.

Xin cảm ơn các thầy cô và bạn học khác từng giúp đỡ và gợi mở cho tôi.

\subsection*{Tài liệu tham khảo}

\begin{enumerate}
  \item Jia Zhipeng, \emph{Sàng tuyến tính và hàm nhân tính}, trao đổi học
  viên WC2012.
  \item Ji Ruyi, \url{http://jiruyi910387714.i11r.com/posts/195270.html},
  jiry\_2's Blog.
  \item Project Euler forum, Problem 521.
\end{enumerate}

\end{document}
